\section{Introduction}

Modern automatic speech recognition (ASR) systems are separated into two
deployment paradigms: cloud-based models that leverage server-scale compute and
edge models that run locally on resource-constrained devices. While cloud ASR
can achieve excellent accuracy by utilizing large models and extensive
computational resources, edge ASR is essential for applications where network
connectivity is unreliable or unavailable, such as offline voice assistants,
medical dictation in remote settings, real-time captioning for accessibility,
and privacy-sensitive voice commands on mobile devices. Edge deployment also
eliminates network round-trip latency and reduces privacy concerns by keeping
audio data on-device.

In edge use cases, latency and transcription quality are the two key---and often
competing---constraints. Achieving human-perceivable real-time performance
requires minimizing time-to-first-token (TTFT) and maintaining low per-token
latency, while simultaneously delivering word error rates (WERs) competitive
with cloud-based alternatives. Balancing these competing objectives on devices
with limited memory, compute, and power budgets remains a central challenge in
practical ASR deployment.

Existing edge ASR models leverage a full-attention encoder architecture, which
allows every frame to directly attend to every other frame in a sequence of
speech audio. This enables powerful contextual disambiguation as it resolves
locally ambiguous acoustics using distant lexical information that occurs
earlier or later in a chunk of speech audio. However, full attention also
introduces quadratic complexity in sequence length and imposes an inherent
``encode-the-whole-utterance'' latency profile: in streaming scenarios, the
encoder must process the entire prefix (or wait for the complete utterance)
before decoder tokens can be emitted, resulting in high TTFT that scales
linearly with utterance length. In practical applications, this reduces system
responsiveness and limits interactivity.

In this paper, we introduce Moonshine v2, a family of ergodic streaming encoder
ASR models designed specifically for latency-critical edge applications.
Moonshine v2 models employ sliding-window attention in a position-free encoder
to enable low-latency streaming inference while maintaining state-of-the-art
accuracy on standard benchmarks. We train three variants of increasing size---
tiny, small, and medium---and show that the models achieve transcription quality and
speed on-par with models 6x their size while running significantly faster (i.e., Whisper Large v3). 
We release the models under a permissive license, encouraging community adoption for on-device, 
latency-critical ASR applications.

The paper is structured as follows. Section~\ref{sec:motivation} analyzes the
latency-accuracy trade-offs inherent in full-attention encoders and motivates
our sliding window approach. Section~\ref{sec:approach} details the Moonshine v2
architecture, including the audio preprocessor, sliding-window encoder, adapter,
and decoder components. Section~\ref{sec:evaluation} presents our experimental
setup and benchmark results across standard ASR datasets. Finally,
Section~\ref{sec:conclusion} discusses implications and future directions for
ergodic streaming ASR.
